%------------------------------------------------------------------------------------
\newpage
\begin{center}
  \textbf{\large ЗАКЛЮЧЕНИЕ}
\end{center}
\refstepcounter{chapter}
\addcontentsline{toc}{chapter}{ЗАКЛЮЧЕНИЕ}
%------------------------------------------------------------------------------------

В ходе выполненной работы была исследована проблема влияния доплеровского масштабирования спектра
на характеристики OFDM сигналов в системах космической связи,
работающих в условиях канала прямой видимости без многолучевого распространения.

Основные результаты работы:
\begin{enumerate}
\item Проведен анализ литературы, посвященной влиянию эффекта Доплера на системы связи,
      в результате которого была выделена и сформулирована задача исследования изолированного
      влияния масштабирования спектра на модуляцию OFDM.
\item Получены выражения
      и оценки уровня интерференции между поднесущими,
      возникающей исключительно из-за масштабирования спектра
      в высокоскоростном LOS-канале с аддитивным белым гауссовским шумом.
\item Методом компьютерного моделирования исследовано непосредственное
      влияние данного эффекта на вероятность битовой ошибки.
      Результаты продемонстрировали деградацию характеристик OFDM системы 
      с ростом относительной скорости и числа поднесущих.
\item Установлены количественные граничные параметры системы,
      при превышении которых влияние доплеровского масштабирования
      спектра становится значимым и должно учитываться при проектировании систем. Существенный вклад в деградацию BER вносит относительная скорость, превышающая $V \approx 1000$ м/с, и количество поднесущих $N > 512$. При меньших значениях этих параметров влиянием эффекта можно пренебречь.
\end{enumerate}

Таким образом, исследование подтвердило, что в высокоскоростных спутниковых каналах связи эффект доплеровского масштабирования спектра является существенным фактором, приводящим к возникновению ICI и снижению помехоустойчивости систем с OFDM модуляцией. Уровень вызванной им интерференции зависит исключительно от относительной скорости движения и структурных параметров сигнала (числа поднесущих), что отличает его от классического доплеровского сдвига.

%Практическая значимость работы заключается в полученных конкретных оценках деградации качества связи, которые могут быть использованы при проектировании и планировании параметров OFDM-систем для группировок спутников, обеспечивающих глобальное покрытие.